\section*{Bab \ref{ch:b3:banach} --- Ruang Banach dan
Teorema Fundamentalnya}

\begin{solution}{exo:b3:banach:1}
(a) Berlaku $\norm{Sx}_2 = \norm x_2$: jadi $S$ merupakan isometri dan $\vertiii S =
1$. Untuk pergeseran mundurnya: $\norm{S^*x}_2^2 = \sum_{n \geq
2}\abs{x_n}^2 \leq \norm x_2^2$, dengan kesamaan pada $x = e_2$:
sehingga $\vertiii{S^*} = 1$.

(b) $\norm{M_ax}_2^2 = \sum\abs{a_n}^2\abs{x_n}^2 \leq \norm
a_\infty^2\norm x_2^2$; dan mengujinya pada $x = e_n$ memberikan $\vertiii{M_a}
\geq \abs{a_n}$ untuk setiap $n$: jadi $\vertiii{M_a} = \norm a_\infty$.

(c) $\abs{\Lambda(f)} \leq \int_0^1\abs f \leq \norm f_\infty$:
sehingga $\norm\Lambda \leq 1$. Untuk $\varepsilon > 0$ misalkan $f_\varepsilon$
bernilai $1$ pada $[0, \frac12 - \varepsilon]$, $-1$ pada $[\frac12 +
\varepsilon, 1]$, dan afin di antaranya: maka $\norm{f_\varepsilon}_\infty =
1$ dan $\Lambda(f_\varepsilon) \geq 1 - 2\varepsilon$: jadi
$\norm\Lambda = 1$. Tidak tercapai: sebab $\Lambda(f) = 1$ dengan $\norm
f_\infty \leq 1$ memaksa $\int_0^{1/2}f = \frac12$ dan
$\int_{1/2}^1 f = -\frac12$, yakni (menurut \omterm{def:b3:topology:continuity}{kekontinuannya} dan $\abs f \leq 1$)
$f \equiv 1$ pada $[0, \frac12]$ dan $f \equiv -1$ pada $[\frac12,
1]$: dan itu kontradiksi di $\frac12$.
\end{solution}

\begin{solution}{exo:b3:banach:2}
Tulis $S = T\bigl(I - T^{-1}(T - S)\bigr)$ dengan
$\vertiii{T^{-1}(T-S)} \leq \vertiii{T^{-1}}\,\vertiii{T - S} =
\theta < 1$: maka menurut \cref{prop:b3:banach:neumann}, $S$ terbalikkan
dengan $S^{-1} = \sum_{n\geq0}\bigl(T^{-1}(T -
S)\bigr)^nT^{-1}$, sehingga
\[
\vertiii{S^{-1} - T^{-1}} \leq
\sum_{n\geq1}\theta^n\,\vertiii{T^{-1}}
= \frac{\vertiii{T^{-1}}^2\,\vertiii{T - S}}{1 - \theta}.
\]
Untuk sistem linearnya: $x_T = T^{-1}b$ dan $x_S = S^{-1}b$
berbeda paling banyak sebesar batas itu dikalikan $\norm b$ --- jadi usikan
kecil pada sistem yang terbalikkan tetap \omterm{def:b3:groups:derived}{terselesaikan} secara tunggal,
dengan kebergantungan Lipschitz penyelesaiannya pada operatornya.
\end{solution}

\begin{solution}{exo:b3:banach:3}
(a) Untuk $\ell^1$: misalkan $(x^{(k)})$ Cauchy. Setiap koordinatnya
Cauchy (sebab $\abs{x^{(k)}_n - x^{(l)}_n} \leq \norm{x^{(k)} -
x^{(l)}}_1$): misalkan $x_n = \lim_kx^{(k)}_n$. Diberikan $\varepsilon$,
untuk $k, l \geq K$: $\sum_{n \leq N}\abs{x^{(k)}_n - x^{(l)}_n}
\leq \varepsilon$ untuk setiap $N$; biarkan $l \to \infty$, lalu $N \to
\infty$: maka $\norm{x^{(k)} - x}_1 \leq \varepsilon$, dan $x =
x^{(k)} - (x^{(k)} - x) \in \ell^1$. Untuk $\ell^\infty$: Cauchy terhadap
$\norm\cdot_\infty$ berarti Cauchy seragam, sehingga konvergen seragam ke
sebuah barisan terbatas. Sedangkan $c_0$ tertutup di $\ell^\infty$: sebab jika
$x^{(k)} \to x$ seragam dengan $x^{(k)}_n \to_n 0$, maka
$\abs{x_n} \leq \norm{x - x^{(k)}}_\infty + \abs{x^{(k)}_n}$
memberikan $\limsup_n\abs{x_n} \leq \varepsilon$: jadi $x \in c_0$; dan
subruang tertutup dari ruang Banach juga Banach. Untuk barisan berhingga:
\omterm{def:b3:topology:interior}{penutupnya} memuat setiap $x \in c_0$ (sebab pemenggalannya konvergen:
$\sup_{n>N}\abs{x_n} \to 0$) dan termuat di $c_0$ yang tertutup.

(b) Lewat kehomogenannya andaikan $\norm x_p = 1$: maka $\abs{x_n} \leq
1$ untuk setiap $n$, sehingga $\abs{x_n}^q \leq \abs{x_n}^p$ dan $\norm
x_q \leq 1 = \norm x_p$; sedangkan untuk $q = \infty$, $\abs{x_n} \leq
\norm x_p$ secara langsung. Kesejatiannya: $x_n = n^{-\alpha}$ dengan
$\frac1q < \alpha \leq \frac1p$ terletak di $\ell^q \setminus
\ell^p$ (lewat deret Riemann).
\end{solution}

\begin{solution}{exo:b3:banach:4}
($\leq$) Jika $\norm f \leq 1$ dan $f\restriction_F = 0$: maka untuk
setiap $y \in F$ berlaku $\abs{f(x)} = \abs{f(x - y)} \leq \norm{x -
y}$; lalu ambil infimumnya. ($\geq$, dan tercapai)
\cref{cor:b3:banach:hbnormed}(3) menghasilkan $f$ dengan
$f\restriction_F = 0$, $\norm f \leq 1$, dan $f(x) = d(x, F)$: sehingga
supremumnya merupakan maksimum. Akibatnya: $F \subseteq \bigcap\{\ker
f : f\restriction_F = 0\}$ secara trivial, dan titik $x \notin F$
tersingkir dari irisannya oleh fungsional di atas
(sebab $f(x) = d(x,F) > 0$, karena $F$ tertutup).
\end{solution}

\begin{solution}{exo:b3:banach:5}
Untuk $y \in \ell^\infty$: $\abs{\Lambda_y(x)} =
\abs{\sum x_ny_n} \leq \norm y_\infty\norm x_1$, sehingga
$\norm{\Lambda_y} \leq \norm y_\infty$; lalu mengujinya pada $e_n$:
$\abs{y_n} = \abs{\Lambda_y(e_n)} \leq \norm{\Lambda_y}$: jadi
kesamaan. Sebaliknya, diberikan $\Lambda \in (\ell^1)'$, tetapkan $y_n =
\Lambda(e_n)$: maka $\abs{y_n} \leq \norm\Lambda$, sehingga $y \in
\ell^\infty$; lalu $\Lambda$ dan $\Lambda_y$ bersesuaian pada barisan
berhingga, yang padat di $\ell^1$ (sebab $\norm{x - \sum_{n\leq
N}x_ne_n}_1 = \sum_{n>N}\abs{x_n} \to 0$): jadi $\Lambda =
\Lambda_y$. Pemetaan $y \mapsto \Lambda_y$ bersifat linear, isometrik,
dan surjektif. Untuk $(\ell^\infty)'$ awal yang sama menghasilkan barisan
$y_n = \Lambda(e_n)$, tetapi barisan berhingga \emph{tidak} padat
di $\ell^\infty$
(sebab barisan konstan $\mathbf 1$ berjarak $1$ dari semuanya),
sehingga $\Lambda$ tidak ditentukan oleh $y_n$ --- dan
memang $(\ell^\infty)' \neq \ell^1$
(\cref{pb:b3:banach:1}).
\end{solution}

\begin{solution}{exo:b3:banach:6}
Untuk setiap $x$ yang tetap, $y \mapsto B(x, y)$ bersifat linear \omterm{def:b3:topology:continuity}{kontinu}:
sehingga $\sup_{\norm y \leq 1}\norm{B(x,y)} < \infty$. Jadi keluarga
$\{B(\cdot, y) : \norm y \leq 1\} \subseteq \mathcal L(E, G)$
(yang setiap anggotanya \omterm{def:b3:topology:continuity}{kontinu}, menurut \omterm{def:b3:topology:continuity}{kekontinuannya} terhadap $x$) bersifat terbatas
titik demi titik pada ruang Banach $E$: sehingga Banach--Steinhaus
(\cref{thm:b3:banach:banachsteinhaus}) melahirkan $C$ dengan
$\norm{B(x,y)} \leq C\norm x$ untuk setiap $\norm y \leq 1$;
lalu kehomogenannya terhadap $y$ menuntaskannya: $\norm{B(x,y)} \leq
C\norm x\norm y$.
\end{solution}

\begin{solution}{exo:b3:banach:7}
(a) Berlaku $\norm f_1 \leq \norm f_\infty$: sehingga identitasnya terbatas
dan bijektif. Inversnya tak terbatas: sebab $f_n(x) = x^n$ mempunyai
$\norm{f_n}_1 = \frac1{n+1} \to 0$ tetapi $\norm{f_n}_\infty = 1$.
Tak ada kontradiksi dengan \cref{cor:b3:banach:equivnorms}: sebab $(\mathcal
C(\intcc01), \norm\cdot_1)$ \emph{tidak} \omterm{def:b3:complete:complete}{lengkap}
(\cref{exo:b3:complete:1}); dan akibat itu menuntut \omterm{def:b3:complete:complete}{kelengkapan}
pada kedua sisinya.

(b) Pada ruang $E_0$ berisi barisan berhingga (yang padat di $\ell^2$),
$\varphi(x) = \sum_n n\,x_n$ bersifat linear dan tak terbatas
(sebab $\varphi(e_n) = n$ dengan $\norm{e_n}_2 = 1$). Teorema grafik
tertutup tak berlaku: sebab $E_0$ tidak \omterm{def:b3:complete:complete}{lengkap} --- dan $\varphi$
tak punya perluasan \omterm{def:b3:topology:continuity}{kontinu} ke $\ell^2$, sehingga menggambarkan bahwa
kepadatan tanpa \omterm{def:b3:topology:continuity}{kekontinuan} seragam tidaklah berdaya
(\cref{thm:b3:complete:extension}).
\end{solution}

\begin{solution}{exo:b3:banach:8}
Kita periksa hipotesis grafik tertutupnya. Misalkan $x_k \to x$ dan
$Tx_k \to z$ di $\ell^2$. Untuk setiap $y$:
\[
\langle z, y\rangle = \lim_k\langle Tx_k, y\rangle
= \lim_k \langle x_k, Ty\rangle = \langle x, Ty\rangle
= \langle Tx, y\rangle,
\]
dengan memakai \omterm{def:b3:topology:continuity}{kekontinuan} hasil kali dalam pada setiap lubangnya
(lewat Cauchy--Schwarz) dan kesimetriannya dua kali. Jadi $z - Tx$ ortogonal
terhadap setiap $y$, khususnya terhadap dirinya sendiri: sehingga $z = Tx$. Grafiknya
tertutup dan $\ell^2$ bersifat Banach: jadi $T$ terbatas
(\cref{thm:b3:banach:closedgraph}). Karena itu operator simetris yang
terdefinisi pada \emph{seluruh} $\ell^2$ otomatis terbatas;
sedangkan operator simetris yang sungguh tak terbatas (posisi, momentum,
Hamiltonian) harus hidup pada subruang padat yang sejati.
\end{solution}

\begin{solution}{exo:b3:banach:9}
Pertama, $\norm{\Lambda_n} = \sum_i\abs{w_{i,n}}$: sebab $\leq$ berasal dari
ketaksamaan segitiga; sedangkan $\geq$ lewat menguji sebuah $f$ linear sepotong-sepotong
dengan $\norm f_\infty \leq 1$ dan $f(x_{i,n}) =
\operatorname{sign}(w_{i,n})$ (interpolasikan secara linear di antara
simpul yang berhingga banyak itu; dan bila simpulnya berimpit tandanya bersesuaian).

($\Rightarrow$) Kekonvergenan titik demi titik pada setiap $f$ mengakibatkan (i),
sekaligus keterbatasan titik demi titik, sehingga Banach--Steinhaus
(\cref{thm:b3:banach:banachsteinhaus}) pada $\mathcal
C(\intcc01)$ yang Banach memberikan (ii).

($\Leftarrow$) Misalkan $M = \sup_n\norm{\Lambda_n} + 1$. Diberikan $f$
dan $\varepsilon$, pilihlah polinomial $P$ dengan $\norm{f -
P}_\infty < \varepsilon/(2M)$
(\cref{cor:b3:complete:weierstrass}); maka
\[
\Bigl|\Lambda_n f - \int_0^1 f\Bigr|
\leq \abs{\Lambda_n(f - P)} + \Bigl|\Lambda_nP -
\int_0^1P\Bigr| + \Bigl|\int_0^1(P - f)\Bigr|
\leq \varepsilon + \Bigl|\Lambda_nP - \int_0^1P\Bigr| \to
\varepsilon .
\]
Untuk bobot positif yang eksak pada konstanta: $\sum_i\abs{w_{i,n}}
= \sum_iw_{i,n} = \Lambda_n(\mathbf 1) = \int_0^1 1 = 1$ bagi
aturan yang eksak pada konstanta --- jadi (ii) berlaku dengan konstanta $1$.
\end{solution}

\begin{solution}{exo:b3:banach:10}
Kelinearan $J$ bersifat formal; sedangkan $\norm{J(x)} = \sup_{\norm f \leq
1}\abs{f(x)} = \norm x$ menurut
\cref{cor:b3:banach:hbnormed}(2). Jika $\dim E = n$: maka $\dim E' =
n$ (sebab sebuah basis memberikan fungsional koordinatnya), sehingga $\dim E'' = n$,
dan $J$ yang injektif (isometrik) menjadi surjektif. Keseparabelannya: misalkan
$(f_n)$ padat di $E'$ lalu pilihlah $\norm{x_n} = 1$ dengan
$\abs{f_n(x_n)} \geq \frac12\norm{f_n}$. Misalkan $F =
\overline{\operatorname{Vect}}(x_n)$; jika $F \neq E$, ambillah $g
\in E'$ dengan $g \neq 0$ yang lenyap pada $F$
(\cref{cor:b3:banach:hbnormed}(3)); lalu pilihlah $f_{n_k} \to g$:
\[
\norm{f_{n_k} - g} \geq \abs{(f_{n_k} - g)(x_{n_k})}
= \abs{f_{n_k}(x_{n_k})} \geq \tfrac12\norm{f_{n_k}}
\geq \tfrac12\bigl(\norm g - \norm{f_{n_k} - g}\bigr),
\]
sehingga $\norm{f_{n_k} - g} \geq \frac13\norm g > 0$: kontradiksi.
Karena itu $F = E$, dan kombinasi rasional (atau $\Q + \iu\Q$) atas
$x_n$ membentuk \omterm{def:b3:topology:interior}{himpunan padat} yang terbilang.
\end{solution}

\begin{solution}{exo:b3:banach:11}
(a) Terdefinisi dengan baik: sebab $d(x, F)$ hanya bergantung pada $\bar x$
(menggeser $x$ oleh $F$ tak mengubah jaraknya).
Kehomogenan dan ketaksamaan segitiganya beralih dari $\norm\cdot$
melalui infimumnya. Pemisahannya memerlukan ketertutupan: sebab $\norm{\bar
x} = 0$ berarti $d(x, F) = 0$, yakni $x \in \bar F = F$, yakni
$\bar x = 0$. Lalu $\vertiii\pi \leq 1$: sebab $\norm{\bar x} \leq
\norm x$. Bola terbuka pada bola terbuka: jika $\norm{\bar x} < 1$,
maka suatu wakilnya mempunyai $\norm{x - y} < 1$; sebaliknya
$\pi(B_E(0,1)) \subseteq B_{E/F}(0,1)$ menurut ketaksamaan
normanya --- jadi $\pi$ terbuka, yakni kasus panutan bagi teorema
pemetaan terbuka.

(b) Misalkan $\sum_k\norm{\bar x_k} < \infty$ lalu angkatlah dengan
$\norm{x_k} \leq \norm{\bar x_k} + 2^{-k}$: maka
$\sum\norm{x_k} < \infty$, sehingga $s = \sum_kx_k$ konvergen di
$E$ yang Banach (\cref{exo:b3:complete:1}(b)), dan \omterm{def:b3:topology:continuity}{kekontinuan}
$\pi$ memberikan $\sum_k\bar x_k = \bar s$: jadi setiap deret $E/F$ yang
konvergen mutlak akan konvergen, dan itu setara dengan
\omterm{def:b3:complete:complete}{kelengkapannya} (lihat latihan yang sama).

(c) Pemetaan $\lambda(x) = \lim_nx_n$ bersifat linear $c \to K$ dan
lenyap tepat pada $c_0$, sehingga ia menginduksikan bijeksi linear
$c/c_0 \to K$. Keisometriannya: $d(x, c_0) = \abs{\lambda(x)}$ ---
untuk $\leq$: kurangkan dari $x$ barisan $x - \lambda(x)\mathbf
1 \in c_0$, sehingga tersisa $\lambda(x)\mathbf 1$ yang bernorma
$\abs{\lambda(x)}$; untuk $\geq$: untuk $y \in c_0$, $\norm{x -
y}_\infty \geq \limsup_n\abs{x_n - y_n} = \abs{\lambda(x)}$.
\end{solution}

\begin{solution}{exo:b3:banach:12}
(a) Kernel $W = \ker P$ tertutup (sebagai prapeta $0$ di bawah
\omterm{def:b3:topology:continuity}{pemetaan kontinu}); sedangkan $V = \operatorname{im}P = \ker(I - P)$
(sebab $Px = x$ jika dan hanya jika $x \in \operatorname{im}P$, dengan memakai $P^2 =
P$), yang juga tertutup. Setiap $x$ terpecah sebagai $Px + (x - Px)$
dengan $Px \in V$ dan $x - Px \in W$, dan $V \cap W = 0$ (sebab $x = Px
= 0$): jadi $E = V \oplus W$.

(b) Argumen grafiknya: misalkan $x_n \to x$ dan $Px_n \to z$.
Maka $z \in V$ (sebab $V$ tertutup dan $Px_n \in V$) dan $x_n - Px_n \to
x - z \in W$ (sebab $W$ tertutup). Jadi $x = z + (x - z)$ dengan $z \in
V$ dan $x - z \in W$; lalu menurut ketunggalan penguraiannya, $z =
Px$. Grafik $P$ tertutup dan $E$ bersifat Banach: sehingga $P$
terbatas (\cref{thm:b3:banach:closedgraph}).

(c) Bagian (a) dan (b) bersama-sama merupakan kesetaraannya. Di ruang Hilbert
setiap $V$ yang tertutup mempunyai komplemen tertutup $V^\perp$
(\cref{ch:b3:hilbert}): jadi setiap subruang tertutup bersifat
berkomplemen. Pada ruang Banach umum hal ini gagal --- sebab $c_0$
tak punya komplemen tertutup di $\ell^\infty$ (yakni teorema
Phillips, yang melampaui perkakas kita) --- sehingga proyeksi terbatas merupakan
hak istimewa, dan teorema grafik tertutup tepat merupakan
pembukuan yang mengubah pemecahan geometri menjadi operator yang terbatas.
\end{solution}

\begin{solution}{pb:b3:banach:1}
\textbf{1.} Untuk $a, b > 0$: menurut kekonkafan $\log$ berlaku
$\log\bigl(\tfrac{a^p}p + \tfrac{b^q}q\bigr) \geq \tfrac1p\log
a^p + \tfrac1q\log b^q = \log(ab)$; lalu eksponenkan. (Jika $ab = 0$
ketaksamaannya trivial.) Kesamaannya terjadi jika dan hanya jika $a^p = b^q$.

\textbf{2.} Kita boleh menganggap $\norm x_p = \norm y_q = 1$
(lewat kehomogenannya; kasus nolnya trivial). Maka
\[
\abs{\langle x, y\rangle} \leq \sum_n\abs{x_n}\abs{y_n}
\leq \sum_n\Bigl(\frac{\abs{x_n}^p}p +
\frac{\abs{y_n}^q}q\Bigr) = \frac1p + \frac1q = 1 =
\norm x_p\norm y_q .
\]
Kesamaannya menuntut $\abs{x_n}^p = \abs{y_n}^q$ untuk setiap $n$
(yakni kasus kesamaan Young) beserta kesejajaran fasa
$x_ny_n$.

\textbf{3.} $\abs{x_n + y_n}^p \leq \bigl(\abs{x_n} +
\abs{y_n}\bigr)\abs{x_n + y_n}^{p-1}$; lalu menjumlahkannya dan menerapkan
Hölder ($p$ terhadap $q$, dengan mencatat $(p - 1)q = p$):
\[
\norm{x + y}_p^p \leq \bigl(\norm x_p + \norm
y_p\bigr)\,\Bigl(\sum_n\abs{x_n + y_n}^{p}\Bigr)^{1/q}
= \bigl(\norm x_p + \norm y_p\bigr)\,\norm{x+y}_p^{p/q};
\]
jika $\norm{x + y}_p \neq 0$ (kalau nol maka trivial), bagilah dengan
$\norm{x+y}_p^{p/q}$ lalu pakai $p - \frac pq = 1$. (Keberhinggaan
$\norm{x+y}_p$ lebih dahulu: $\abs{x_n+y_n}^p \leq
2^p(\abs{x_n}^p + \abs{y_n}^p)$.)

\textbf{4.} \omterm{def:b3:complete:complete}{Kelengkapannya}: seperti untuk $\ell^1$
(\cref{exo:b3:banach:3}), lewat limit koordinat demi koordinat ditambah
batas ekor seragam $\sum_{n\leq N}\abs{x^{(k)}_n - x^{(l)}_n}^p
\leq \varepsilon^p$, dengan membiarkan $l$ lalu $N$ menuju tak hingga.
Kepadatan barisan berhingga: $\norm{x - \sum_{n\leq
N}x_ne_n}_p^p = \sum_{n > N}\abs{x_n}^p \to 0$.

\textbf{5.} Hölder memberikan $\abs{\Lambda_y(x)} \leq \norm
x_p\norm y_q$: jadi $\norm{\Lambda_y} \leq \norm y_q$. Untuk mengujinya: misalkan
$x^{(N)}_n = \abs{y_n}^{q-1}\overline{\operatorname{sign}}(y_n)$
untuk $n \leq N$ dan $0$ selebihnya (dengan $\operatorname{sign}$ berupa
fasa unimodularnya, sehingga $x_ny_n = \abs{y_n}^q$). Maka
$\Lambda_y(x^{(N)}) = \sum_{n\leq N}\abs{y_n}^q$ dan
$\norm{x^{(N)}}_p = \bigl(\sum_{n\leq
N}\abs{y_n}^{q}\bigr)^{1/p}$ (sebab $(q-1)p = q$), sehingga
\[
\norm{\Lambda_y} \geq \Bigl(\sum_{n\leq
N}\abs{y_n}^q\Bigr)^{1 - 1/p} \xrightarrow[N\to\infty]{}
\norm y_q .
\]

\textbf{6.} Tetapkan $y_n = \Lambda(e_n)$. Dengan vektor uji yang
sama, $\sum_{n \leq N}\abs{y_n}^q = \Lambda(x^{(N)}) \leq
\norm\Lambda\,\bigl(\sum_{n\leq N}\abs{y_n}^q\bigr)^{1/p}$,
sehingga $\bigl(\sum_{n\leq N}\abs{y_n}^q\bigr)^{1/q} \leq
\norm\Lambda$ untuk setiap $N$: jadi $y \in \ell^q$ dengan $\norm y_q \leq
\norm\Lambda$. Fungsional $\Lambda$ dan $\Lambda_y$ bersesuaian
pada barisan berhingga yang padat (pertanyaan 4): sehingga $\Lambda =
\Lambda_y$. Bersama pertanyaan 5, $y \mapsto \Lambda_y$ merupakan
isomorfisma isometrik $\ell^q \cong (\ell^p)'$.

\textbf{7.} Misalkan $\xi \in (\ell^p)''$. Dengan mengomposisikannya dengan
isometri $\ell^q \cong (\ell^p)'$ pada pertanyaan 6, $\xi$ mendefinisikan
sebuah unsur $(\ell^q)'$, yang (menurut pertanyaan 6 dengan $p, q$
ditukar) berupa $\Lambda_z$ untuk $z \in \ell^p$ yang tunggal: jadi untuk setiap
$y \in \ell^q$, $\xi(\Lambda_y) = \sum_nz_ny_n$. Di pihak
lain $J(z)(\Lambda_y) = \Lambda_y(z) = \sum_ny_nz_n$: yaitu nilai yang
sama. Karena setiap unsur $(\ell^p)'$ berupa suatu $\Lambda_y$, maka
$\xi = J(z)$: jadi $J$ surjektif --- sehingga $\ell^p$ \omterm{rem:b3:banach:bidual}{refleksif}.

\textbf{8.} Barisan berhingga yang entrinya rasional (pada bagian real dan
imajinernya) bersifat terbilang dan padat di $\ell^p$ (untuk $p < \infty$) dan
di $c_0$. Di $\ell^\infty$: keluarga $\{\mathbf 1_A : A
\subseteq \N\}$ bersifat tak terbilang dengan $\norm{\mathbf 1_A -
\mathbf 1_B}_\infty = 1$ untuk $A \neq B$; sehingga bola $B(\mathbf
1_A, \frac12)$ saling lepas berpasangan, dan \omterm{def:b3:topology:interior}{himpunan padat} harus memotong
masing-masingnya: jadi tak ada \omterm{def:b3:topology:interior}{himpunan padat} yang terbilang.

\textbf{9.} Seandainya $\ell^1$ \omterm{rem:b3:banach:bidual}{refleksif}, maka $(\ell^\infty)'
\cong ((\ell^1)')' = J(\ell^1)$ akan \omterm{def:b3:galois:separable}{separabel} (sebagai peta
isometrik dari $\ell^1$ yang \omterm{def:b3:galois:separable}{separabel}); lalu menurut
\cref{exo:b3:banach:10}, keseparabelan \omterm{def:b3:banach:operator}{dual}
$(\ell^\infty)'$ akan memaksa $\ell^\infty$ \omterm{def:b3:galois:separable}{separabel} ---
dan itu bertentangan dengan pertanyaan 8. Jadi $\ell^1$ tidak \omterm{rem:b3:banach:bidual}{refleksif} (dan
$(\ell^\infty)'$ sungguh lebih besar daripada $\ell^1$, seperti dikonkretkan
pertanyaan 11).

\textbf{10.} Kehomogenan $p$ jelas; sedangkan keaditifan bawahnya:
rata-rata bersifat linear, dan $\limsup(u_n + v_n) \leq \limsup u_n +
\limsup v_n$. Pada $c$: rata-rata Cesàro sebuah barisan konvergen
konvergen ke limitnya, sehingga $p(x) = \lim x =
\mathrm{LIM}(x)$ di sana; khususnya $\mathrm{LIM} \leq p$ pada
$c$. Lalu Hahn--Banach (\cref{thm:b3:banach:hahnbanach}) memperluas
$\mathrm{LIM}$ ke $\ell^\infty_\R$ dengan $\mathrm{LIM} \leq p$
secara global. Kepositifannya: untuk $x \geq 0$, $-\mathrm{LIM}(x) =
\mathrm{LIM}(-x) \leq p(-x) = \limsup\text{rata-rata}(-x) \leq 0$.
Keinvarianan pergeserannya: rata-rata $x - Sx$ berteleskop menjadi
$\frac{x_1 - x_{n+1}}n \to 0$, sehingga $p(\pm(x - Sx)) = 0$ dan
$\mathrm{LIM}(x - Sx) = 0$. Batasnya: $\mathrm{LIM}(x) \leq p(x)
\leq \limsup x$ (sebab rata-rata tertinggal di belakang supremumnya), lalu menerapkan hal ini
pada $-x$ memberikan batas bawahnya.

\textbf{11.} Karena $e_n \in c_0$, maka $\mathrm{LIM}(e_n) = 0$ untuk setiap
$n$. Jika $\mathrm{LIM} = \Lambda_y$ dengan $y \in \ell^1$, maka
$y_n = \Lambda_y(e_n) = 0$ untuk setiap $n$: sehingga $\Lambda_y = 0$; padahal
$\mathrm{LIM}(\mathbf 1) = 1$. Jadi $\mathrm{LIM} \in
(\ell^\infty)' \setminus \{\Lambda_y : y \in \ell^1\}$:
sekali lagi $(\ell^\infty)' \neq \ell^1$.

\textbf{12.} Untuk $x = (0,1,0,1,\dots)$: $x + Sx = \mathbf 1$,
sehingga $2\,\mathrm{LIM}(x) = \mathrm{LIM}(x) + \mathrm{LIM}(Sx) =
1$: jadi $\mathrm{LIM}(x) = \frac12$. Seandainya $\varphi$ sebuah
perluasan limit yang multiplikatif dan invarian terhadap pergeseran: maka $x\cdot
Sx = 0$ memberikan $\varphi(x)\varphi(Sx) = \varphi(x)^2 = 0$, sehingga
$\varphi(x) = 0$; padahal $\varphi(x) + \varphi(Sx) =
\varphi(\mathbf 1) = 1$ memberikan $2\varphi(x) = 1$: kontradiksi.
Jadi perata-rataan dan perkalian tak dapat hidup berdampingan.

\textbf{13.} Misalkan $\norm{x^{(k)}}_p \leq M$. Koordinatnya
terbatas oleh $M$: sehingga ekstraksi diagonal memberikan subbarisan
(yang masih ditulis $x^{(k)}$) dengan $x^{(k)}_n \to x_n$ untuk setiap
$n$. Lalu $x \in \ell^p$: sebab $\sum_{n\leq N}\abs{x_n}^p =
\lim_k\sum_{n\leq N}\abs{x^{(k)}_n}^p \leq M^p$ untuk setiap $N$.
Kekonvergenan lemahnya: untuk $y \in \ell^q$ dan $N$ sembarang,
\[
\abs{\Lambda_y(x^{(k)} - x)} \leq
\Bigl|\sum_{n \leq N}(x^{(k)}_n - x_n)y_n\Bigr|
+ 2M\Bigl(\sum_{n>N}\abs{y_n}^q\Bigr)^{1/q},
\]
dengan suku pertamanya menuju $0$ ketika $k \to \infty$ (sebab koordinatnya
berhingga banyak) dan suku keduanya kecil untuk $N$ besar
(lewat Hölder pada ekornya): jadi $\Lambda_y(x^{(k)}) \to \Lambda_y(x)$
untuk setiap $y$ --- yakni kekonvergenan lemah, sebab setiap fungsionalnya berupa
$\Lambda_y$ (pertanyaan 6). Di $\ell^1$ hal ini gagal: tinjaulah
$(e_n)$ yang terbatas. Setiap subbarisan $(e_{n_k})$ konvergen
koordinat demi koordinat ke $0$, sehingga satu-satunya calon limit lemahnya adalah $0$;
padahal menguji terhadap $y \in \ell^\infty$ yang didefinisikan oleh $y_{n_k} =
(-1)^k$ (dan $0$ selebihnya) memberikan $\Lambda_y(e_{n_k}) = (-1)^k$
yang divergen. Jadi tak ada subbarisan yang konvergen lemah: sehingga \omterm{thm:b3:topology:metriccompact}{kekompakan
barisan} secara lemah pada bolanya mencirikan dunia yang \omterm{rem:b3:banach:bidual}{refleksif}.

\textbf{14.} Untuk $y \in \ell^1$, $\Lambda_y(x) = \sum x_ny_n$
terdefinisi pada $c_0$ dengan $\abs{\Lambda_y(x)} \leq \norm
x_\infty\norm y_1$, dan mengujinya pada $x^{(N)} =
(\operatorname{sign}y_1, \dots, \operatorname{sign}y_N, 0,
\dots) \in c_0$ memberikan $\Lambda_y(x^{(N)}) =
\sum_{n\leq N}\abs{y_n} \to \norm y_1$: jadi
$\norm{\Lambda_y} = \norm y_1$. Sebaliknya, untuk $\Lambda \in
(c_0)'$ tetapkan $y_n = \Lambda(e_n)$; maka uji yang sama memberikan
$\sum_{n \leq N}\abs{y_n} = \Lambda(x^{(N)}) \leq
\norm\Lambda$, sehingga $y \in \ell^1$, dan $\Lambda = \Lambda_y$ pada
barisan berhingga yang padat, jadi di mana-mana. Rantai
\omterm{def:b3:banach:operator}{dualnya}: $(c_0)' = \ell^1$, $(\ell^1)' = \ell^\infty$
(\cref{exo:b3:banach:10}), dan $(\ell^\infty)' \supsetneq \ell^1$
(pertanyaan 9--11): jadi kerefleksifannya gagal sejak langkah pertama
--- sebab $c_0'' = \ell^\infty \neq c_0$ --- dan tak pernah pulih.

\textbf{15.} Keluarga $J x^{(k)} \in E'' = (E')'$ berisi
fungsional terbatas pada ruang \emph{Banach} $E'$ (sebab \omterm{def:b3:banach:operator}{dual}
selalu \omterm{def:b3:complete:complete}{lengkap}, \cref{prop:b3:banach:llcomplete}); dan untuk setiap
$\Lambda \in E'$, barisan $Jx^{(k)}(\Lambda) =
\Lambda(x^{(k)})$ konvergen, sehingga terbatas.
Lalu Banach--Steinhaus pada $E'$ memberikan $\sup_k\norm{Jx^{(k)}} <
\infty$, dan $J$ bersifat isometrik
(\cref{rem:b3:banach:bidual}): jadi $\sup_k\norm{x^{(k)}} <
\infty$.

\textbf{16.} ($\Rightarrow$) Keterbatasannya adalah pertanyaan 15;
sedangkan koordinatnya berupa fungsional $\Lambda_{e_n}$.
($\Leftarrow$) Misalkan $M = \sup_k\norm{x^{(k)}}_p$, $y \in
\ell^q$, dan $\varepsilon > 0$; pilihlah $N$ dengan
$\bigl(\sum_{n>N}\abs{y_n}^q\bigr)^{1/q} < \varepsilon$. Maka
\[
\abs{\Lambda_y(x^{(k)} - x)} \leq
\sum_{n\leq N}\abs{x^{(k)}_n - x_n}\,\abs{y_n} +
\norm{x^{(k)} - x}_p\,\varepsilon
\leq \sum_{n\leq N}\abs{x^{(k)}_n - x_n}\abs{y_n} +
(M + \norm x_p)\,\varepsilon,
\]
dan jumlah berhingganya menuju $0$: sehingga $\limsup \leq (M + \norm
x_p)\varepsilon$ untuk setiap $\varepsilon$. (Bahwa $x \in
\ell^p$ dengan $\norm x_p \leq M$ menyusul dari batas
penggalan berhingga bergaya Fatou: $\sum_{n \leq N}\abs{x_n}^p =
\lim_k\sum_{n\leq N}\abs{x^{(k)}_n}^p \leq M^p$.) Untuk $e_k$
di $\ell^2$: terbatas dan koordinatnya $\to 0$, sehingga $e_k
\rightharpoonup 0$, padahal $\norm{e_k} = 1$: jadi satu satuan massanya
lolos ke indeks tak berhingga, tak terlihat oleh fungsional
mana pun yang tetap.

\textbf{17.} Ambil $\Lambda$ dengan $\norm\Lambda = 1$ dan
$\Lambda(x) = \norm x$ (\cref{cor:b3:banach:hbnormed}). Maka
$\norm x = \lim\Lambda(x^{(k)}) \leq
\liminf\norm\Lambda\,\norm{x^{(k)}} =
\liminf\norm{x^{(k)}}$. (Dengan $e_k \rightharpoonup 0$: $0
\leq \liminf 1$, dan ketaksamaannya dapat sejati.)

\textbf{18.} Di $\ell^2$, $\norm{x^{(k)} - x}_2^2 =
\norm{x^{(k)}}_2^2 - 2\operatorname{Re}\langle x^{(k)},
x\rangle + \norm x_2^2$ (pada kasus realnya: $-2\langle x^{(k)},
x\rangle$). Kekonvergenan lemah yang diterapkan pada fungsional
$\Lambda_x$ memberikan $\langle x^{(k)}, x\rangle \to \norm
x_2^2$, sedangkan normanya konvergen menurut hipotesis: sehingga ruas kanannya
menuju $\norm x^2 - 2\norm x^2 + \norm x^2 = 0$.
Contoh penyangkal tanpa kekonvergenan normanya: $e_k
\rightharpoonup 0$ dengan $\norm{e_k - 0} = 1 \not\to 0$.

\textbf{19.} Kekonvergenan koordinatnya: terapkan fungsional
$\Lambda_{e_n} \in (\ell^1)' = \ell^\infty$ (dengan $y = e_n$).
Konstruksinya: setelah $k_{j-1}, N_{j-1}$ dipilih, ambillah $k_j >
k_{j-1}$ sedemikian besar sehingga $\sum_{n \leq
N_{j-1}}\abs{x^{(k_j)}_n} < \frac\delta{20}$ (sebab koordinatnya
berhingga banyak dan masing-masing $\to 0$), lalu $N_j > N_{j-1}$ sedemikian besar
sehingga ekornya memenuhi $\sum_{n > N_j}\abs{x^{(k_j)}_n} <
\frac\delta{20}$ (menurut kekonvergenan deret yang mendefinisikan
$\norm{x^{(k_j)}}_1$). Blok $B_j = \intoc{N_{j-1}}{N_j}$
lalu membawa segalanya kecuali $\frac\delta{10}$ dari massa $x^{(k_j)}$.

\textbf{20.} Dengan $y$ seperti didefinisikan (sehingga $\abs{y_n} \leq 1$
di mana-mana):
\[
\langle x^{(k_j)}, y\rangle
= \sum_{n \in B_j}\abs{x^{(k_j)}_n}
+ \sum_{n \notin B_j}x^{(k_j)}_ny_n
\geq \Bigl(\norm{x^{(k_j)}}_1 - \frac\delta{10}\Bigr) -
\frac\delta{10}
\geq \delta - \frac{2\delta}{10} = \frac{4\delta}5 > 0,
\]
dengan ketaksamaan di tengahnya karena massa di luar bloknya paling
banyak $\frac\delta{10}$ (pertanyaan 19). Padahal $y \in
\ell^\infty = (\ell^1)'$ dan $x^{(k)} \rightharpoonup 0$
memaksa $\langle x^{(k_j)}, y\rangle \to 0$: kontradiksi.
Karena itu barisan $\ell^1$ yang nol secara lemah juga nol dalam norma, dan lewat
translasi barisan yang konvergen lemah akan konvergen dalam norma: itulah teorema
Schur.

\textbf{21.} Subbarisan $(e_k)$ yang konvergen lemah akan
berlimit $0$ (menurut koordinatnya), sehingga menurut Schur $\norm{e_{k_j}}_1
\to 0$ --- padahal normanya $1$. Jadi tak ada subbarisan yang konvergen
lemah, seperti ditemukan dengan tangan pada pertanyaan 13. Tak ada
paradoks: Schur mengatakan bahwa barisan tak dapat membedakan \omterm{def:b3:topology:topology}{topologi} lemah
dari \omterm{def:b3:topology:topology}{topologi} norma di $\ell^1$ (sedangkan \omterm{def:b3:topology:topology}{topologinya} sendiri
memang berbeda --- \omterm{def:b3:topology:topology}{persekitaran} lemahnya tak pernah terbatas normanya), dan
\omterm{thm:b3:topology:metriccompact}{kekompakan barisan} secara lemah pada bola satuannya merupakan sifat lain
yang lebih kuat, yang setara dengan kerefleksifan
(Eberlein--\v Smulian, yang melampaui perkakas kita; kegagalannya,
setidaknya, sudah kita buktikan).

\textbf{22.} \emph{Sensusnya.}
\begin{center}
\begin{tabular}{l|l|c|c|c}
$E$ & $E'$ & sep. & refl. & bola \omterm{def:b3:topology:compact}{kompak} barisan lemah\\
\hline
$c_0$ & $\ell^1$ & ya & tidak & tidak ($e_1{+}\dots{+}e_k$)\\
$\ell^1$ & $\ell^\infty$ & ya & tidak & tidak ($e_k$, p.~21)\\
$\ell^p$ & $\ell^q$ & ya & ya & ya (p.~13)\\
$\ell^\infty$ & $\supsetneq\ell^1$ & tidak & tidak & tidak\\
\end{tabular}
\end{center}
Ciri khasnya: $c_0$ --- bidualnya $\ell^\infty$, yakni langkah
tak \omterm{rem:b3:banach:bidual}{refleksif} yang pertama (pertanyaan 14); $\ell^1$ --- sifat
Schur (pertanyaan 20); $\ell^p$ --- kerefleksifan dan kekompakan
lemah (pertanyaan 7 dan 13); $\ell^\infty$ ---
ketakseparabelan dan limit Banach: yakni fungsional yang tak dapat diwakili
barisan mana pun (pertanyaan 8 dan 10--11). Satu keluarga ruang,
empat dunia yang berbeda.

\textbf{23.} Misalkan $(f_k) \subseteq F$ dengan $\norm{x - f_k}
\to d$. Maka $\norm{f_k} \leq \norm x + \sup_k\norm{x -
f_k}$: jadi terbatas, sehingga menurut pertanyaan 13 ada subbarisan $f_{k_j}
\rightharpoonup f$. Jika $f \notin F$, maka $\delta =
\operatorname{dist}(f, F) > 0$ (sebab $F$ tertutup); lalu pada $F \oplus \R
f$ bentuk linear $\lambda(g + tf) = t$ memenuhi
$\abs{\lambda(u)} \leq \norm u/\delta$ (sebab $\norm{g +
tf} \geq \abs t\,\delta$), dan Hahn--Banach memperluasnya menjadi
$\Lambda \in (\ell^p)'$ dengan $\Lambda\restriction_F = 0$ dan
$\Lambda(f) = 1$; tetapi lalu $0 = \Lambda(f_{k_j}) \to
\Lambda(f) = 1$: kontradiksi. Jadi $f \in F$, dan $x -
f_{k_j} \rightharpoonup x - f$ memberikan, menurut pertanyaan 17,
\[
d \leq \norm{x - f} \leq \liminf_j\,\norm{x - f_{k_j}} = d :
\]
jadi jaraknya tercapai di $f$. Di $c_0$: $\abs{\Lambda(x)}
\leq \sum_n2^{-n}\abs{x_n} < \norm x_\infty$ untuk setiap $x
\neq 0$ (sebab barisan nol yang tak nol tak mungkin memenuhi $\abs{x_n} =
\norm x_\infty$ untuk setiap $n$), sedangkan pemenggalan $(1,
\dots, 1, 0, \dots)$ memberikan $\Lambda = 1 - 2^{-N} \to 1$: jadi
$\norm\Lambda = 1$, yang tak pernah tercapai. Rumus jaraknya: untuk $f
\in \ker\Lambda$, $\abs{\Lambda(x)} = \abs{\Lambda(x - f)}
\leq \norm{x - f}$, sehingga $\operatorname{dist}(x, \ker\Lambda)
\geq \abs{\Lambda(x)}$; sebaliknya, untuk $u$ di bola satuannya
dengan $\Lambda(u) \geq 1 - \varepsilon$, vektor $f = x -
\frac{\Lambda(x)}{\Lambda(u)}u$ terletak di $\ker\Lambda$ dengan
$\norm{x - f} \leq \frac{\abs{\Lambda(x)}}{1 - \varepsilon}$:
jadi kesamaan. Seandainya suatu $f \in \ker\Lambda$ mencapainya, maka $z = x -
f$ akan memenuhi $\abs{\Lambda(z)} = \abs{\Lambda(x)} =
\norm z \neq 0$, sehingga $\Lambda$ akan mencapai normanya di
$z/\norm z$: mustahil. Jadi ada hiperbidang tertutup $c_0$ yang tak punya
titik terdekat di mana pun --- sehingga kerefleksifan bukanlah hiasan.

\textbf{24.} Tetapkan $y_1 = x^{(1)}$. Diberikan $y_1, \dots, y_j$,
setiap pemetaan $k \mapsto \langle y_i, x^{(k)}\rangle$ menuju $0$
(sebab $y_i \in \ell^2 = (\ell^2)'$), sehingga ada $k_{j+1}$ melewati
indeks sebelumnya dengan $\abs{\langle y_i,
x^{(k_{j+1})}\rangle} \leq \frac1{j+1}$ untuk $i = 1, \dots,
j$; sebutlah pilihan itu $y_{j+1}$. Maka
\[
\Bigl\lVert\sum_{j=1}^my_j\Bigr\rVert_2^2
= \sum_{j=1}^m\norm{y_j}_2^2
+ 2\sum_{j=2}^m\sum_{i<j}\langle y_i, y_j\rangle
\leq mC^2 + 2\sum_{j=2}^m\frac{j-1}j
\leq mC^2 + 2m,
\]
lalu dengan membaginya dengan $m^2$: $\norm{\frac1m\sum_jy_j}_2^2 \leq
\frac{C^2 + 2}m \to 0$. (Untuk limit lemah $x \neq 0$, terapkan
hal ini pada $x^{(k)} - x$.) Pada $(e_k)$ yang ortonormal bahkan tak perlu
ekstraksi: $\norm{\frac1m(e_1 + \dots +
e_m)}_2 = \frac{\sqrt m}m = \frac1{\sqrt m}$. Jadi rata-ratanya
mengubah kekonvergenan lemah menjadi kekonvergenan norma: itulah
sifat Banach--Saks pada $\ell^2$.

\textbf{25.} Karena $A_mx = \frac1m\sum_{i=0}^{m-1}S^ix$, maka
kelinearan dan keinvarianan pergeserannya memberikan $L(A_mx) = L(x)$. Untuk sembarang
$u$ yang terbatas dan $\varepsilon > 0$, pilihlah $N$ dengan $u_n \leq
\limsup u + \varepsilon$ untuk $n \geq N$; lalu kepositifan yang diterapkan
pada $(\limsup u + \varepsilon)\mathbf 1 - S^Nu \geq 0$ beserta
$L(S^Nu) = L(u)$ memberikan $L(u) \leq \limsup u + \varepsilon$,
dan secara simetris $L(u) \geq \liminf u - \varepsilon$: sehingga
$\liminf A_mx \leq L(x) \leq \limsup A_mx$ untuk setiap $m$. Jika
$x$ berperiode $T$, maka $A_Tx$ berupa barisan konstan yang sama dengan
rata-rata periodenya $\mu$: jadi $L(x) = \mu$ untuk \emph{setiap} limit
Banach --- yakni $\frac13$ pada $(1,0,0,\dots)$ dan $\frac12$ pada
$(0,1,0,1,\dots)$ seperti pada pertanyaan 12. Untuk barisan bloknya:
di $N = 3^j$ dengan $j$ genap blok terakhirnya seluruhnya satu, sehingga
rata-rata Cesàronya $\geq \frac{3^j - 3^{j-1}}{3^j} = \frac23$;
sedangkan di $N = 3^j$ dengan $j$ ganjil semua angka satunya duduk di $\intoc0
{3^{j-1}}$, sehingga rata-ratanya $\leq \frac13$. Karena itu $p(x) \geq
\frac23$ dan $-p(-x) = \liminf_n\frac{x_1 + \dots + x_n}n
\leq \frac13$. Pada $M = c \oplus \R x$ definisikan $\Lambda_+(y +
tx) = \lim y + t\,p(x)$. Dominasinya oleh $p$: untuk $t > 0$,
kesublinearannya memberikan $p(tx) \leq p(y + tx) + p(-y)$, yakni
$p(y + tx) \geq t\,p(x) + \lim y$ (perhatikan $p(\pm y) = \pm\lim
y$ untuk $y \in c$: sebab rata-rata Cesàro barisan konvergen
konvergen ke limitnya); untuk $t = -s < 0$, $p(y) \leq p(y -
sx) + p(sx)$ memberikan $p(y - sx) \geq \lim y - s\,p(x)$; sedangkan untuk $t
= 0$ terjadi kesamaan. Jadi $\Lambda_+ \leq p$ pada $M$, dan
Hahn--Banach memperluasnya menjadi $L_+ \leq p$ pada $\ell^\infty_\R$,
yang merupakan limit Banach persis seperti pada pertanyaan 10 (sebab dominasi
oleh $p$ melahirkan kepositifan, keinvarianan pergeseran, dan nilai
$\lim$ pada $c$), dengan $L_+(x) = p(x) \geq \frac23$. Perhitungan yang sama
dengan $\Lambda_-(y + tx) = \lim y - t\,p(-x)$
(memakai $p(-sx) \leq p(y - sx) + p(-y)$ untuk kasus $t = -s
< 0$) melahirkan limit Banach $L_-$ dengan $L_-(x) = -p(-x) \leq
\frac13$. Dua limit Banach, satu barisan, dua nilai:
jadi di luar dunia periodik (dan, lebih umum, dunia yang \emph{hampir
konvergen}), sebuah limit Banach merupakan pilihan yang sesungguhnya.

\end{solution}
